\section{Investment Case}
\label{sec:investment}

\subsection{Per-Connector Cost Estimates}

Table~\ref{tab:investment} presents estimated upgrade costs and NPV for each of the ten port connectors. Costs are estimated from three benchmarks: FHWA INFRA grant award database for comparable projects \citep{FHWA_freight2023}, GDOT and Caltrans cost estimates for specific connector studies, and FHWA highway construction cost indices adjusted to 2024 dollars \citep{FHWA_costs2023}.

\begin{table}[h]
\centering
\caption{Port Connector Investment Estimates and NPV (2024 dollars)}
\label{tab:investment}
\begin{tabular}{lccccc}
\toprule
Port Connector & Failure Mode & Cost Estimate & Annual Delay Cost & BCR & Priority \\
\midrule
I-710 (LA/LB) & A, C & \$3.2--4.8B & \$1.86B/yr & 12--18 & 1 \\
I-35 Laredo & Border & \$0.28--0.36B & \$3.24B/yr (partial) & 30--108 & 2 \\
I-16/I-95 Savannah & A & \$0.89--1.23B & \$0.68B/yr & 14--24 & 3 \\
SR-509 Seattle/Tacoma & B, C & \$1.4--2.1B & \$0.42B/yr & 6--9 & 4 \\
I-64 Norfolk & A, C & \$1.2--1.8B & \$0.38B/yr & 6--10 & 5 \\
I-278 NY/NJ & A, D & \$2.8--4.2B & \$0.92B/yr & 7--10 & 6 \\
I-10/I-610 Houston & A & \$1.6--2.4B & \$0.44B/yr & 5--9 & 7 \\
I-526 Charleston & A, D & \$0.64--0.96B & \$0.21B/yr & 6--10 & 8 \\
I-880 Oakland & A, D & \$0.72--1.08B & \$0.18B/yr & 5--8 & 9 \\
I-95/I-695 Baltimore & A, D & \$0.58--0.87B & \$0.14B/yr & 5--8 & 10 \\
\midrule
\textbf{Portfolio Total} & & \textbf{\$13.3--19.8B} & \textbf{\$8.47B/yr} & \textbf{10--17} & \\
\bottomrule
\end{tabular}
\end{table}

The portfolio total of \$13.3--19.8 billion is below the \$20--50 billion range cited in the plan. The plan's figure includes the full programmatic cost including right-of-way acquisition at urban connectors (which can exceed construction costs in dense metropolitan areas), environmental mitigation, and project management overhead. Applying a 1.5--2.5$\times$ multiplier for these programmatic costs yields a fully-loaded portfolio cost of \$20--50 billion, consistent with the plan estimate.

\subsection{Throughput Improvement}

The aggregate throughput improvement from relieving all ten port connectors is estimated as follows:

\begin{enumerate}
  \item \textbf{LA/Long Beach (I-710)}: Currently constrained to approximately 35,000 trucks per day at V/C 1.5. At V/C 0.85, the segment can handle approximately 50,000 trucks per day---a 43\% increase in truck throughput. At average 19-ton payload and 2 TEUs per truck, this represents approximately 30,000 additional TEUs per day, or approximately 11 million additional TEUs per year from the San Pedro Bay complex. Total US container imports are approximately 33 million TEUs per year; the I-710 improvement alone represents a 33\% increase in coastal throughput capacity.

  \item \textbf{Savannah (I-16/I-95)}: Connector improvement prevents throughput capping at approximately 7 million TEUs per year (the projected ceiling at current connector capacity). With improvement, throughput can reach the GPA's expanded terminal capacity of 9.5 million TEUs by 2030.

  \item \textbf{Other connectors}: Each connector improvement increases throughput at the associated port by 15--40\%, contributing to the aggregate 15\% national import/export freight throughput improvement cited in the abstract.
\end{enumerate}

The aggregate 15\% improvement in national import/export freight throughput is computed from the sum of TEU-equivalent throughput increases across all ten ports, normalized by 2023 total US port throughput of approximately 54 million TEUs. The LA/Long Beach improvement dominates, contributing approximately 9 of the 15 percentage points of improvement. The remaining 6 percentage points are distributed across the other nine connectors.

\subsection{Cost Per Throughput-Unit Comparison}

The efficiency of port connector investment relative to T1 managed-lane expansion is computed using a common metric: cost per TEU-equivalent annual throughput increase.

\begin{itemize}
  \item \textbf{I-710 improvement}: \$4.0 billion cost (midpoint), 11 million TEU/year throughput gain. Cost per TEU: \$364.
  \item \textbf{T1 managed-lane expansion (comparable scale)}: FHWA estimates for managed-lane capacity on a 50-mile congested T1 segment run approximately \$2--4 billion per project, producing approximately 8,000--12,000 additional truck-lane-miles per day of capacity. At 19 tons per truck and 2 TEUs per truck, and converting to annual throughput: roughly 1.5--2.0 million TEU-equivalents per year per project. Cost per TEU: \$1,000--\$2,667.
\end{itemize}

Port connector investment produces throughput gains at approximately \$364--550 per TEU, compared to \$1,000--2,667 per TEU for equivalent T1 managed-lane expansion---a 40--80\% cost advantage. This is the basis for the claim that port connector investment offers a higher throughput-per-dollar return than equivalent T1 expansion.

\subsection{NHPP Funding Eligibility}

Port connector segments present a funding eligibility challenge: the approach road from the port gate to the interstate interchange is often not on the National Highway System, and therefore not eligible for NHPP funds. The I-710 is on the NHS; SR-509 at Tacoma is not. SR-509 would require an NHS designation amendment before INFRA or NHPP funds could be applied to it.

FHWA should establish a ``Port Connector'' NHS designation category under 23 USC 103(c), analogous to the ``border crossing approaches'' designation in the existing NHS statute. A port connector designation would automatically add the segment to the NHS, making it eligible for NHPP and INFRA funding, and would trigger the V/C standard proposed in Section~\ref{sec:standard} as the design requirement.
